%==================== CHAPTER 1 ====================
\ChapterPage{Introduction}

\section{Background}
Even in the most modern and efficient business teams today, there is still a considerable amount of repetitive work that can't be entirely delegated. A recruiting HR manager may review hundreds of resumes for one job and find that if they are doing the same thing each time; it takes hours to do the job. A marketing team writes the same campaign slogan 5 times, each for different channels: LinkedIn, Instagram, Newsletter, etc. A podcast producer puts together an interview with a guest over the course of a few hours by collating interviews, articles and social profiles spread across the web. None of this is very creative, but it is the time people waste on the decisions that really count.

The software these teams were using didn't really help as much. ATS software, HR software or CRM software is great at storing and retrieval of records, not so good at handling unstructured data like a resume, a job description in natural language or a long article. They have rigid rules and will not work with inputs that do not fit its requirements. In recent years, the large language model has evolved, making it possible to achieve much more here: \cite{autogen,guo,react}. Such models can understand unstructured text, think about it, invoke external tools and share its output with other agents in a pipeline. Our project, in collaboration with \Partner, takes this ability to the business processes where the software doesn't match the real requirement of the teams.

\section{Motivation}
The current challenges with tools are no mere nuisance. When it comes to hiring, keyword screening has been found to be useful in rewarding those candidates who know how to craft a resume rather than the best candidates for the position Gan et al. \cite{gan} showed that an LLM-based pipeline can screen resumes eleven times faster than a human and reach a respectable F1 score, but their system was tested only on IT-sector data and could not generate rubrics from a job description or conduct an interview. Wilson and Caliskan \cite{wilson} found that embedding-based screening favoured White-associated names in 85 out of every 100 cases and that Black male candidates were disadvantaged in up to 100\% of test cases. Mu et al. \cite{mu} then showed that adversarial content hidden inside a resume could manipulate LLM screening more than 80\% of the time, and that the best available defences reduced this by only around a quarter. If a recruiter is given a ranked shortlist without a note, there is nothing they can do to detect these mistakes until they impact a real life employee.

In marketing, content that is sent out without a proper review can cause harm to a brand that can take months to fix. Research for podcasts can also result in a reputation risk if the host or podcast relies on unsourced or misattributed claims \cite{marketingfm,cove}. Whether it's the social, or the finance, or the other aspect of the business, the data required to work is in different systems with different logins, API, rate limits, quirks etc. The task of compiling everything in one place is quite a full time occupation.

What teams really need is something that can process the repetitive parts end-to-end, something which could be explained unambiguously to a person for purposes of review and trusted, and then wait for the person to sign off on it before some important thing occurs. This project is meant to fill that void.

\section{Aim and Objectives}
\RunIn{Aim:}
\textbf{\ProjectTitle}'s mission is to create an integrated multi-agent AI platform that handles repetitive, high-volume components of HR recruitment, marketing content production, podcast guest research and enterprise data integration — leaving a human reviewer to handle each decision that matters. The platform is built with a series of specialized LLM agents, approved by the approval gates, and supported by a retrieval-augmented knowledge layer, with audit logs and controls in between.
\vspace{12pt}
\RunIn{Objectives:}
The specific objectives of the project are:
\begin{itemize}
  \item To build HR workflows that generate weighted rubrics from a job description, screen and score resumes with cited evidence, conduct AI-assisted interviews and produce a ranked shortlist for the recruiter to review, alongside assistants that answer onboarding, leave and HR policy questions.
  \item To build marketing workflows that turn a theme or source content into a content calendar, channel-specific drafts and audience segments, with every piece of content held for human approval before it goes anywhere.
  \item To build podcast research workflows that gather public sources about a guest, index them for retrieval, and produce a structured briefing in which every claim is linked to a source and labelled as a verified fact, an allegation or an opinion.
  \item To build a connector layer that communicates with Zoho, Workday-style and Greenhouse-style systems through authenticated, rate-limit-aware, idempotent API calls, with human approval required before any write operation goes through.
  \item To enforce consistent governance across all workflows through a compliance reviewer that checks for PII, prompt injection and unsupported claims, an approval queue that holds every consequential output, and an audit log that records every prompt, tool call, source and reviewer action.
  \item To package the entire platform as a reusable workflow framework and measure its performance on accuracy, hallucination rate, latency, cost and task completion.
\end{itemize}
